\scanpage{52}
Nous allons examiner \deleted{d'ailleurs} ce que donne le théor. de Tate, pour le groupe de Lubin Tate $G_{\unsure{\varepsilon}}$.
Notons que, si $K_\pi\subset\overline K$ est l'extension \inserted{abélienne} de groupe $A^*$ définie par $G$ (via $\mathrm{Gal}(\overline K/K)\xrightarrow{\chi_\pi}A^*$), alors \deleted{\unsure{un}} $T_p(G)(\overline K)$ est un module inversible sur $A$, et les bases de celui-ci sont permutées de façon simplement transitive par $\mathrm{Gal}(K_\pi/K)$. Donc, il y a un choix privilégié (à isom \uline{unique} près) d'une extension algébrique $K_\pi$ de $K$ (indépendant du choix d'une clôture $\overline K$ !) avec un générateur
\[
e\in T_p(G)^*(K_\pi)
\]
[en fait, $\mathrm{Spec}\,K_\pi$ représente le foncteur $T_p(G)^*=\varprojlim G(\nu)^*$, où le $*$ désigne les générateurs \deleted{\unsure{comme}} $A$-modulaires]. Nous supposons fixé $K_\pi$ de cette façon, et choisissons pour $\overline K$ une surextension de $K_\pi$. Donc le choix de $e$ donne un isom de $A$-modules
\[
\overbrace{T_p(G)(K_\pi)}^{T_p}\simeq A,\qquad\lambda\longmapsto\lambda e.
\]
\margin{\diagram{\begin{tikzcd}
T_p\arrow[d,equal]&T_p\otimes A\arrow[d,equal]&T_p\otimes C\arrow[d,equal]\\
A\arrow[r,no head]\arrow[d,no head]&A\otimes_{\mathbb Z_p}A\arrow[r,no head]\arrow[d,no head]&A\otimes_{\mathbb Z_p}C\arrow[d,no head]\\
\mathbb Z_p\arrow[r,no head]&A\arrow[r,no head]&C
\end{tikzcd}}}
On regarde sur $T_p(\unsure{G})\otimes_A C$ d'une part l'action de $A$ provenant de l'action de $A$ sur $G$, (``action à opérateurs'') d'autre part l'action provenant de la structure de $C$-vectoriel. Enfin, on a l'action de $\pi=\mathrm{Gal}(\overline K/K)$, qui commute à l'action (opér.) de $A$. Supposons $K/\mathbb Q_p$ galoisien pour simplifier.
\mspage{51}
